\section{紧嵌入}
\label{sec:sobolev-compact-embeddings}

连续嵌入说明，$W^{1,p}$中的有界集合在目标空间中仍然有界。有界还不足以抽出收敛子列：第~\ref{chap:functions-as-vectors}章已经看到，无限维空间的闭球不紧。本节研究更强的结论：Sobolev范数有界的函数列，能否在较弱的目标范数下选出强收敛子列。先在次临界指数下建立这一结论，再用一个集中序列说明临界指数处它为什么失败，最后把结论转写为整类函数的有限覆盖。

\subsection{\texorpdfstring{$L^q$}{Lq}中的强收敛子列}
\label{sec:sobolev-compact-lq}

导数的统一控制阻止函数在小尺度上剧烈振荡。直观地说，导数有界的函数在小尺度上几乎不变，把它们在固定尺度上平均以后，就只剩有限个自由度；有限维中的有界集可以用有限个点逼近。下面的证明把这个直观分成三步：延拓到全空间，用平移差统一控制磨光误差，再对固定尺度的磨光函数作有限维近似。

\begin{lemma}[\texorpdfstring{$L^p$}{Lp}中的收敛子列]
\label{lem:sobolev-rellich-lp}
设$\Omega\subset\mathbb R^d$为有界Lipschitz区域，$1\leq p<\infty$。则$W^{1,p}(\Omega)$中的每个有界序列都有在$L^p(\Omega)$中收敛的子列。
\end{lemma}

\begin{proof}
设$\lVert f_n\rVert_{W^{1,p}(\Omega)}\leq M$。由第~\ref{chap:metric-spaces}章，完备度量空间中全有界的子集是相对紧的，其中每个序列都有收敛子列；$L^p(\Omega)$完备，所以只需证明$\{f_n\}$在$L^p(\Omega)$中全有界。

\emph{延拓。}令$w_n=Ef_n$，由定理~\ref{thm:sobolev-extension}，$\lVert w_n\rVert_{W^{1,p}(\mathbb R^d)}\leq\lVert E\rVert M=:M'$，且所有$w_n$在同一有界集合$B_\Omega$外为零。限制到$\Omega$不增大$L^p$距离，所以只需证明$\{w_n\}$在$L^p(\mathbb R^d)$中全有界。

\emph{平移差的统一控制。}对$w\in C_c^1(\mathbb R^d)$，$w(x+z)-w(x)=\int_0^1\nabla w(x+tz)\cdot z\,\mathrm dt$。两边取$L^p$范数，把范数移进对$t$的积分，得
\begin{equation}
\lVert w(\cdot+z)-w\rVert_p\leq|z|\,\lVert\nabla w\rVert_p .
\label{eq:sobolev-translation-gradient}
\end{equation}
两边对$W^{1,p}$范数连续，由稠密性，它对每个$w\in W^{1,p}(\mathbb R^d)$成立。磨光误差$w*\rho_\varepsilon-w=\int(w(\cdot-z)-w)\rho_\varepsilon(z)\,\mathrm dz$是平移差的加权平均，所以
\[
\lVert w_n*\rho_\varepsilon-w_n\rVert_p\leq\sup_{|z|\leq\varepsilon}\lVert w_n(\cdot-z)-w_n\rVert_p\leq\varepsilon\lVert\nabla w_n\rVert_p\leq\varepsilon d^{1-1/p}M' .
\]
这一估计对所有$n$一致成立：导数的统一界把所有函数的磨光误差统一控制住了。

\emph{有限维近似。}固定$\varepsilon$，令$v_n=w_n*\rho_\varepsilon$。它们都在同一个有界方体$Q$外为零。由Hölder不等式，
\[
|v_n(x)|\leq\lVert\rho_\varepsilon\rVert_{p'}\lVert w_n\rVert_p\leq\lVert\rho_\varepsilon\rVert_{p'}M',\qquad
|\nabla v_n(x)|\leq\lVert\nabla\rho_\varepsilon\rVert_{p'}M'=:\Lambda ,
\]
这族函数一致有界并有统一的Lipschitz常数$\Lambda$。把$Q$分成边长为$h$的$N$个小方体$Q_l$，取中心$y_l$，令$Pv=\sum_lv(y_l)\mathbf 1_{Q_l}$。在每个$Q_l$上$|v-v(y_l)|\leq\Lambda\sqrt d\,h$，所以$\lVert v_n-Pv_n\rVert_p\leq\Lambda\sqrt d\,h\,|Q|^{1/p}$。而$Pv_n$都落在$N$维空间$\operatorname{span}\{\mathbf 1_{Q_l}\}$中，坐标的绝对值一致有界；有限维空间中的有界集是全有界的。所以给定精度，先取$h$使$\Lambda\sqrt dh|Q|^{1/p}$足够小，再覆盖$\{Pv_n\}$，就得到$\{v_n\}$的有限网。

\emph{合并。}任给$\eta>0$，先取$\varepsilon$使$\varepsilon d^{1-1/p}M'<\eta/2$，再取$\{v_n\}$的有限$\eta/2$网$\{g_1,\ldots,g_m\}$。对每个$n$，选$g_i$与$v_n$相距小于$\eta/2$，则$\lVert w_n-g_i\rVert_p<\eta$。所以$\{w_n\}$全有界。
\end{proof}

证明清楚地显示了各条件的职责。导数的统一界使磨光误差与$n$无关；区域有界使固定尺度的函数族只需有限个小方体描述；目标范数$L^p$比$W^{1,p}$弱，才允许用磨光后的函数近似原函数，而不必近似导数。

\begin{theorem}[紧嵌入]
\label{thm:sobolev-rellich}
设$d\geq2$，$1<p<d$，$1\leq q<p^*=dp/(d-p)$，$\Omega\subset\mathbb R^d$为有界Lipschitz区域。则$W^{1,p}(\Omega)$中的每个有界序列都有在$L^q(\Omega)$中收敛的子列。
\end{theorem}

\begin{proof}
由引理~\ref{lem:sobolev-rellich-lp}，有界序列$(f_n)$有在$L^p$中收敛的子列，仍记为$(f_n)$。若$q\leq p$，由区域有限测度，$\lVert g\rVert_q\leq|\Omega|^{1/q-1/p}\lVert g\rVert_p$，子列在$L^q$中也是Cauchy列。若$p<q<p^*$，取$\theta\in(0,1)$使$1/q=\theta/p+(1-\theta)/p^*$。把$|g|^q=|g|^{\theta q}|g|^{(1-\theta)q}$用于指数$p/(\theta q)$与$p^*/((1-\theta)q)$的Hölder不等式（这两个指数共轭），得
\[
\lVert g\rVert_q\leq\lVert g\rVert_p^\theta\lVert g\rVert_{p^*}^{1-\theta}.
\]
取$g=f_n-f_m$。由定理~\ref{thm:sobolev-embedding-subcritical}，$\lVert f_n-f_m\rVert_{p^*}\leq2CM$有界，而$\lVert f_n-f_m\rVert_p\to0$，$\theta>0$，所以$\lVert f_n-f_m\rVert_q\to0$。$L^q$完备，子列在$L^q$中收敛。
\end{proof}

称这样的包含为\term{紧嵌入}（compact embedding）：包含映射把有界集映为相对紧集，按第~\ref{chap:linear-operators}章的说法，它是紧算子。同样的论证在$p=1$时也成立；$p\geq d$时，由区域有界可把$W^{1,p}$放进每个$W^{1,s}$（$s<d$），于是对每个有限的$q$都有紧嵌入。

紧嵌入只给出目标空间中的强收敛，不给出Sobolev范数下的收敛。取$f_n(x)=n^{-1}\sin(nx_1)$，则$\lVert f_n\rVert_p\leq|\Omega|^{1/p}/n\to0$，$\partial_1f_n=\cos(nx_1)$，其余偏导数为零，序列在$W^{1,p}(\Omega)$中有界。若某个子列在$W^{1,p}$中收敛，极限只能是$L^p$极限$0$，于是$\lVert\cos(nx_1)\rVert_p\to0$。但$|\cos|\leq1$给出$\int_\Omega\cos^2(nx_1)\leq|\Omega|^{1-1/p}\lVert\cos(nx_1)\rVert_p$，而由Riemann--Lebesgue引理，$\int_\Omega\cos^2(nx_1)\,\mathrm dx=\frac12|\Omega|+\frac12\int_\Omega\cos(2nx_1)\,\mathrm dx\to\frac12|\Omega|$，矛盾。

紧嵌入常与弱收敛配合使用。$1<p<\infty$时$W^{1,p}(\Omega)$是自反空间，由第~\ref{chap:linear-functionals}章，有界序列有弱收敛子列；紧嵌入进一步让同一子列在$L^q$中强收敛，并且两个极限相同。上例中$f_n$在$W^{1,p}$中弱收敛到$0$，在$L^q$中强收敛到$0$，在$W^{1,p}$范数下却不收敛。

\subsection{临界指数处的集中现象}
\label{sec:sobolev-concentration}

定理~\ref{thm:sobolev-rellich}要求$q<p^*$，是严格不等式。在临界指数$q=p^*$处，嵌入仍连续，却不再是紧的。反例来自把一个固定形状压缩到一点、同时按恰当比例放大高度。

沿用上一小节的$d\geq2$与$1<p<d$。取$x_0\in\Omega$与$r>0$使$B(x_0,r)\subset\Omega$，取非零$\varphi\in C_c^\infty(\mathbb R^d)$，支撑在$B(0,r)$内。对正整数$n$，令
\[
f_n(x)=n^{(d-p)/p}\varphi\bigl(n(x-x_0)\bigr).
\]
$f_n$支撑在$B(x_0,r/n)\subset\Omega$内。换元$y=n(x-x_0)$，$\mathrm dx=n^{-d}\mathrm dy$，逐项计算：
\begin{equation*}
\begin{aligned}
\lVert f_n\rVert_p&=n^{(d-p)/p}n^{-d/p}\lVert\varphi\rVert_p=n^{-1}\lVert\varphi\rVert_p,\\
\lVert\nabla f_n\rVert_p&=n^{(d-p)/p}\,n\,n^{-d/p}\lVert\nabla\varphi\rVert_p=\lVert\nabla\varphi\rVert_p,
\end{aligned}
\end{equation*}
\[
\lVert f_n\rVert_{p^*}=n^{(d-p)/p}n^{-d/p^*}\lVert\varphi\rVert_{p^*}=\lVert\varphi\rVert_{p^*},
\]
最后一式用了$d/p^*=d/p-1=(d-p)/p$。所以$(f_n)$在$W^{1,p}(\Omega)$中有界，$L^p$范数按$n^{-1}$衰减，梯度的$L^p$范数和函数的$L^{p^*}$范数都保持不变。对每个$x\neq x_0$，$n$充分大时$x$不在支撑内，$f_n(x)=0$，所以$f_n\to0$几乎处处。

若某个子列在$L^{p^*}(\Omega)$中收敛到$g$，由区域有限测度，它也在$L^p$中收敛到$g$；而它在$L^p$中趋于$0$，所以$g=0$，这与$\lVert f_n\rVert_{p^*}=\lVert\varphi\rVert_{p^*}\neq0$矛盾。因此在临界范数下没有强收敛子列。函数的质量向一点集中：在$L^{p^*}$中看起来一直一样大，却不收敛到任何函数。

这说明严格指数条件不是证明技术的产物。在次临界的$q<p^*$下，同一序列满足$\lVert f_n\rVert_q=n^{(d-p)/p-d/q}\lVert\varphi\rVert_q$，指数$(d-p)/p-d/q=d/p^*-d/q<0$，所以$f_n\to0$在$L^q$中成立，与定理~\ref{thm:sobolev-rellich}一致。若去掉区域有界，即使不集中也会失去紧性：在$\mathbb R^d$上把同一个函数沿一个方向平移得越来越远，范数不变，支撑两两不交，任何子列都不是Cauchy列。

\subsection{有界函数类的有限覆盖}
\label{sec:sobolev-finite-covers}

紧嵌入是关于序列的结论，也可以转写为关于整类函数的结论。取定理~\ref{thm:sobolev-rellich}的条件，令$\mathcal F_R=\{f\in W^{1,p}(\Omega):\lVert f\rVert_{W^{1,p}}\leq R\}$，把它看作$L^q(\Omega)$的子集。

$\mathcal F_R$中的每个序列都有在$L^q$中收敛的子列，极限不必属于$\mathcal F_R$。由第~\ref{chap:metric-spaces}章的相对紧性判据，这说明$\mathcal F_R$在$L^q(\Omega)$中相对紧；在完备空间中，相对紧等价于全有界。所以对每个$\varepsilon>0$，覆盖数$N(\varepsilon,\mathcal F_R,L^q)$有限：存在有限个函数，使$\mathcal F_R$中每个函数与其中之一的$L^q$距离不超过$\varepsilon$。

这里得到的只是“数量有限”。覆盖数随$\varepsilon\to0$怎样增长，紧性论证没有给出：引理~\ref{lem:sobolev-rellich-lp}的证明中，磨光半径与网格尺寸都依赖精度，计数时还要把每一步的代价统一起来。要得到增长阶，需要具体构造代表，第~\ref{chap:function-class-covers}章给出了一般的构造方法，第~\ref{chap:besov-spaces}章将在周期Besov球上用小波系数算出对数覆盖数的确切增长阶。

两个看似相近的结论不能代替有限覆盖。临界处的连续嵌入只说明$\mathcal F_R$在$L^{p^*}$中有界，上一小节的集中序列说明它在$L^{p^*}$中不全有界：若全有界，它在完备的$L^{p^*}$中相对紧，集中序列就会有收敛子列。弱紧性也不能代替：$1<p<\infty$时$\mathcal F_R$在$W^{1,p}$中弱紧，但它在$W^{1,p}$范数下不全有界，第一小节的$n^{-1}\sin(nx_1)$（乘以适当常数使其落在$\mathcal F_R$中）在$W^{1,p}$中没有收敛子列，由完备性也就没有Cauchy子列，而全有界集合中的每个序列都有Cauchy子列。
